% ===== PARTIE 1 : objectif et méthodes (Personne A) =====
\section{Objectif}
Une pompe fonctionne pendant 40 minutes et sa puissance a été relevée toutes
les 10 minutes. L'énergie consommée est l'aire sous la courbe de puissance :
\[
  E = \int_0^{T} P(t)\,dt .
\]
Comme on ne dispose que de points mesurés, on estime cette intégrale avec deux
méthodes de quadrature programmées en Scilab : les trapèzes et Simpson.

\section{Méthodes}
\paragraph{Trapèzes.} Sur chaque intervalle $[t_i, t_{i+1}]$, on relie les deux
points par un segment :
\[
  T = \sum_{i=1}^{n} (t_{i+1}-t_i)\,\frac{y_i+y_{i+1}}{2}.
\]
Cette méthode accepte des pas inégaux.

\paragraph{Simpson.} On remplace la courbe par des paraboles sur des paires
d'intervalles :
\[
  S = \frac{h}{3}\Bigl[y_1 + y_{n+1} + 4(y_2+y_4+\dots+y_n)
      + 2(y_3+y_5+\dots+y_{n-1})\Bigr].
\]
Elle exige des pas égaux et un nombre $n$ pair de sous-intervalles. La fonction
programmée refuse les autres cas avec un message d'erreur.
